% Reconciled reader additions. Source witnesses remain untouched.
\clearpage
\begingroup
\tagtrue{FOL}
\begin{editorial}
आगे प्रमाण सिद्धांत का स्रोत में उपलब्ध प्रारूप पूरा दिया गया है। स्रोत की
संपादकीय चेतावनियाँ मूल परियोजना की स्थिति बताती हैं; मुख्य पीडीएफ़ से बाहर
होने का उनका उल्लेख इस हिंदी संस्करण पर लागू नहीं होता। प्रारूप में चिह्नित
अपूर्णता और प्रयोगात्मकता को अनुवाद ने समाप्त मान नहीं लिया है।
\end{editorial}
\olimport[content/proof-theory]{proof-theory}
\olimport[content/proof-theory/proof-search]{proof-search}
\olchapter{hir}{ptr}{प्रमाण सिद्धांत के अतिरिक्त स्रोत-पाठ}
\olfileid[hi]{pt}{cut}{int-alt}
\olsection{कट-विलोपन के प्रमाण का अतिरिक्त अंश}
\begin{editorial}
यह जमी हुई स्रोत प्रति में अलग रखा गया प्रमाण-अंश है। इसे उसी रूप में
सुरक्षित रखा गया है; यह अपने-आप में एक नया पूर्ण प्रमेय या पूरा प्रमाण नहीं है।
\end{editorial}
\subimport{content/proof-theory/cut-elimination/}{intuitionistic.tex}
\olimport[content/proof-theory/propositions-as-types]{sequent-natural-deduction}
\olfileid[hi]{pt}{rules-extra}{intro}
\olsection{अनुक्रम कलन की अतिरिक्त नियम-सारणियाँ}
\begingroup
\olfileid[hi]{pt}{rules-extra}{G1i}
\olimport[content/proof-theory/sequent-calculus]{rules-G1i}
\clearpage
\endgroup
\begingroup
\olfileid[hi]{pt}{rules-extra}{G3i}
\olimport[content/proof-theory/sequent-calculus]{rules-G3i}
\clearpage
\endgroup
\begingroup
\olfileid[hi]{pt}{rules-extra}{LK}
\olimport[content/proof-theory/sequent-calculus]{rules-LK}
\clearpage
\endgroup
\begingroup
\olfileid[hi]{pt}{rules-extra}{mG3i}
\olimport[content/proof-theory/sequent-calculus]{rules-mG3i}
\clearpage
\endgroup
\endgroup
\olpart{hir}{पूरक और वैकल्पिक स्रोत-पाठ}
\begin{editorial}
इस भाग में वे भिन्न पाठ हैं जो जमी हुई स्रोत प्रति में मौजूद हैं, किंतु उसके
सामान्य संकलन-क्रम में नहीं आते। कहीं विषय पहले के अध्याय से मिलता है, पर
प्रमाण, उदाहरण या प्रस्तुति भिन्न है; इसलिए उसे बिना परिवर्तन सुरक्षित रखा गया है।
\end{editorial}
\begingroup
\tagtrue{FOL}
\olchapter{hir}{sfr}{समुच्चयों, फलनों और आगमन पर पूरक पाठ}
\olimport[content/sets-functions-relations/sets]{proofs-about-sets}
\olimport[content/sets-functions-relations/functions]{isomorphic-functions}
\olimport[content/sets-functions-relations/inductive-defs-proofs]{introduction}
\endgroup
\begingroup
\tagfalse{FOL}
\olchapter{hir}{pl}{प्रतिज्ञप्ति तर्कशास्त्र: निर्दोषता और पूर्णता के वैकल्पिक पाठ}
\olimport[content/propositional-logic/syntax-and-semantics]{soundness}
\olimport[content/propositional-logic/syntax-and-semantics]{completeness}
\endgroup
\begingroup
\tagtrue{FOL}
\olchapter{hir}{fol}{प्रथम-क्रम तर्कशास्त्र के पूरक पाठ}
\olimport[content/first-order-logic/syntax-and-semantics]{introduction}
\olimport[content/first-order-logic/axiomatic-deduction]{provability}
\olimport[content/first-order-logic/completeness]{maximally-consistent-sets}
\endgroup
\begingroup
\tagtrue{FOL}
\olchapter{hir}{mod}{अंकगणित के गैर-मानक निदर्श: वैकल्पिक प्रस्तुति}
\olimport[content/model-theory/basics]{nonstandard-arithmetic}
\endgroup
\begingroup
\tagtrue{FOL}
\olchapter{hir}{inc}{निरूपणीयता: फलनों का वर्ग C}
\olimport[content/incompleteness/representability-in-q]{c}
\olimport[content/incompleteness/representability-in-q]{c-representable}
\endgroup
\begingroup
\tagtrue{FOL}
\olchapter{hir}{sol}{द्वितीय-क्रम भाषा: पूरक परिचय}
\olimport[content/second-order-logic/syntax-and-semantics]{language-of-sol}
\endgroup
\begingroup
\tagtrue{FOL}
\olchapter{hir}{lam}{लैम्ब्डा कलन के पूरक पाठ}
\olimport[content/lambda-calculus/syntax]{conversion}
\olimport[content/lambda-calculus/lambda-definability]{lists}
\endgroup
\begingroup
\tagtrue{FOL}
\olchapter{hir}{mvl}{बहुमूल्य तर्कशास्त्र: प्रमाण-सैद्धांतिक धारणाएँ}
\olimport[content/many-valued-logic/sequent-calculus]{proof-theoretic-notions}
\endgroup
\begingroup
\tagtrue{FOL}
\olchapter{hir}{nml}{सामान्य मोडल तर्कशास्त्र: पूरक पाठ}
\olimport[content/normal-modal-logic/syntax-and-semantics]{normal-modal-logics}
\endgroup
\begingroup
\tagtrue{FOL}
\olchapter{hir}{int}{अंतःप्रज्ञावादी अर्थविज्ञान: प्रतिज्ञप्तियाँ}
\olimport[content/intuitionistic-logic/semantics]{propositions}
\endgroup
\olchapter{hir}{routes}{वैकल्पिक अध्यापन-क्रम और संपादकीय टिप्पणियाँ}
अगली टिप्पणियाँ स्रोत के वैकल्पिक अध्यापन-क्रमों से हैं। उनके साझा खंड इस
पुस्तक में पहले ही दिए गए हैं; केवल अध्याय-चालक बदलने के कारण उन्हें फिर से
नहीं दोहराया गया है।
\section{समुच्चय, संबंध और फलन: प्रारंभिक क्रम}
\begingroup
\RenewDocumentCommand\olimport{s o m o}{}
\RenewDocumentCommand\olpart{m m}{}
\RenewDocumentCommand\olchapter{m m m}{}
\def\OLEndChapterHook{}\def\OLEndPartHook{}
\subfile{content/sets-functions-relations/sets-functions-relations.tex}
\endgroup
\section{संबंधों का अध्याय-क्रम}
\begingroup
\RenewDocumentCommand\olimport{s o m o}{}
\RenewDocumentCommand\olpart{m m}{}
\RenewDocumentCommand\olchapter{m m m}{}
\def\OLEndChapterHook{}\def\OLEndPartHook{}
\subfile{content/sets-functions-relations/relations/relations.tex}
\endgroup
\section{समुच्चयों का आकार: प्रारंभिक क्रम}
\begingroup
\RenewDocumentCommand\olimport{s o m o}{}
\RenewDocumentCommand\olpart{m m}{}
\RenewDocumentCommand\olchapter{m m m}{}
\def\OLEndChapterHook{}\def\OLEndPartHook{}
\subfile{content/sets-functions-relations/size-of-sets/size-of-sets.tex}
\endgroup
\section{वाक्यविन्यास और अर्थविज्ञान का संयुक्त क्रम}
\begingroup
\RenewDocumentCommand\olimport{s o m o}{}
\RenewDocumentCommand\olpart{m m}{}
\RenewDocumentCommand\olchapter{m m m}{}
\def\OLEndChapterHook{}\def\OLEndPartHook{}
\subfile{content/first-order-logic/syntax-and-semantics/syntax-and-semantics.tex}
\endgroup
